% GENERATED FILE. Edit canonical lessons, then run npm run generate:books.
% canonical-source-hash: fnv1a64:9ef3c41e842780f6
% canonical-lessons: AR-C338-tabin, AR-C338-darih, AR-C338-wafah, AR-C338-maqbara, AR-C338-nash

\chapter{Remembering the dead}
\label{ch:ar-remembering-the-dead}
\hlchaptermodality{338}

\begin{chapteropening}
\textbf{By the end of this chapter:} \emph{I can say a eulogy, a tomb, a death, a cemetery and a coffin.}

Complete the last lesson of chapter 338: I can say a eulogy, a tomb, a death, a cemetery and a coffin.
\end{chapteropening}

\section[\texorpdfstring{\ar{تأبين} (taʾbīn)}{تأبين (taʾbīn)}]{\ar{تأبين} (taʾbīn) — a eulogy, a memorial tribute}

\label{lesson:AR-C338-tabin}



\begin{quote}
Before the new one: say the Arabic for an heir, then the Arabic for a widow.
\end{quote}

\subsection*{You'll want to know: \ar{تأبين}}
\textbf{\ar{تأبين}} — \emph{taʾbīn} — \textquotedblleft{}a eulogy, a memorial tribute\textquotedblright{}.

\begin{quote}
\ar{حضر} \ar{الكتاب} \ar{حفل} \ar{تأبين} \ar{الشاعر}.
\end{quote}

\begin{quote}
\emph{ḥaḍara al-kuttāb ḥafl taʾbīn ash-shāʿir.}
\end{quote}

\begin{quote}
The writers attended the poet’s memorial.
\end{quote}

\begin{cousinweb}
\textbf{\ar{تأبين}} (\emph{taʾbīn}) is built on the native Arabic root \emph{ʾ-b-n} (Lane), where Form II \emph{abbana} is \textquotedblleft{}to praise a person after his death\textquotedblright{}; \emph{taʾbīn} is its verbal noun.
\end{cousinweb}

\begin{grammarlens}[title={What you've built}]
One more word for telling what happened.
\end{grammarlens}

\subsection*{Your turn}
\begin{itemize}
\raggedright
  \item \emph{Say it:} \emph{taʾbīn}
  \item \emph{Say it:} \emph{taʾbīn}, once more
  \item \emph{Say it:} \emph{ḥaḍara al-kuttāb ḥafl taʾbīn ash-shāʿir.} — The writers attended the poet’s memorial.
  \item \emph{Recall:} say the Arabic for an heir, then the Arabic for a widow, then say \emph{taʾbīn} again
\end{itemize}

\begin{tcolorbox}[breakable,colback=teal!4,colframe=teal!35!black,title={Before you move on}]
What does \ar{تأبين} mean? (\textquotedblleft{}A eulogy, a memorial tribute\textquotedblright{}.) Say it once more.
\end{tcolorbox}
\section[\texorpdfstring{\ar{ضريح} (ḍarīḥ)}{ضريح (ḍarīḥ)}]{\ar{ضريح} (ḍarīḥ) — a tomb, a mausoleum}

\label{lesson:AR-C338-darih}



\begin{quote}
Before the new one: say the Arabic for a widow, then the Arabic for a eulogy.
\end{quote}

\subsection*{You'll want to know: \ar{ضريح}}
\textbf{\ar{ضريح}} — \emph{ḍarīḥ} — \textquotedblleft{}a tomb, a mausoleum\textquotedblright{}.

\begin{quote}
\ar{زار} \ar{السياح} \ar{ضريح} \ar{الملك} \ar{القديم}.
\end{quote}

\begin{quote}
\emph{zāra as-suyyāḥ ḍarīḥ al-malik al-qadīm.}
\end{quote}

\begin{quote}
The tourists visited the old king’s tomb.
\end{quote}

\begin{cousinweb}
\textbf{\ar{ضريح}} (\emph{ḍarīḥ}) is built on the native Arabic root \emph{ḍ-r-ḥ} (Lane), \textquotedblleft{}to dig a trench for the dead\textquotedblright{}; a \emph{ḍarīḥ} is first that grave, then a built tomb.
\end{cousinweb}

\begin{grammarlens}[title={What you've built}]
One more word for telling what happened.
\end{grammarlens}

\subsection*{Your turn}
\begin{itemize}
\raggedright
  \item \emph{Say it:} \emph{ḍarīḥ}
  \item \emph{Say it:} \emph{ḍarīḥ}, once more
  \item \emph{Say it:} \emph{zāra as-suyyāḥ ḍarīḥ al-malik al-qadīm.} — The tourists visited the old king’s tomb.
  \item \emph{Recall:} say the Arabic for a widow, then the Arabic for a eulogy, then say \emph{ḍarīḥ} again
\end{itemize}

\begin{tcolorbox}[breakable,colback=teal!4,colframe=teal!35!black,title={Before you move on}]
What does \ar{ضريح} mean? (\textquotedblleft{}A tomb, a mausoleum\textquotedblright{}.) Say it once more.
\end{tcolorbox}
\section[\texorpdfstring{\ar{وفاة} (wafāh)}{وفاة (wafāh)}]{\ar{وفاة} (wafāh) — a death, a passing}

\label{lesson:AR-C338-wafah}



\begin{quote}
Before the new one: say the Arabic for a eulogy, then the Arabic for a tomb.
\end{quote}

\subsection*{You'll want to know: \ar{وفاة}}
\textbf{\ar{وفاة}} — \emph{wafāh} — \textquotedblleft{}a death, a passing\textquotedblright{}.

\begin{quote}
\ar{بكى} \ar{الجميع} \ar{بعد} \ar{وفاة} \ar{الجدة}.
\end{quote}

\begin{quote}
\emph{bakā al-jamīʿ baʿda wafāt al-jadda.}
\end{quote}

\begin{quote}
Everyone wept after the grandmother’s death.
\end{quote}

The polite noun beside \ar{موت}.

\begin{cousinweb}
\textbf{\ar{وفاة}} (\emph{wafāh}) is built on the native Arabic root \emph{w-f-y} (Lane), \textquotedblleft{}to be complete, to fulfil\textquotedblright{}. \emph{Wafāh} is a life brought to its full term.
\end{cousinweb}

\begin{grammarlens}[title={What you've built}]
One more word for telling what happened.
\end{grammarlens}

\subsection*{Your turn}
\begin{itemize}
\raggedright
  \item \emph{Say it:} \emph{wafāh}
  \item \emph{Say it:} \emph{wafāh}, once more
  \item \emph{Say it:} \emph{bakā al-jamīʿ baʿda wafāt al-jadda.} — Everyone wept after the grandmother’s death.
  \item \emph{Recall:} say the Arabic for a eulogy, then the Arabic for a tomb, then say \emph{wafāh} again
\end{itemize}

\begin{tcolorbox}[breakable,colback=teal!4,colframe=teal!35!black,title={Before you move on}]
What does \ar{وفاة} mean? (\textquotedblleft{}A death, a passing\textquotedblright{}.) Say it once more.
\end{tcolorbox}
\section[\texorpdfstring{\ar{مقبرة} (maqbara)}{مقبرة (maqbara)}]{\ar{مقبرة} (maqbara) — a cemetery}

\label{lesson:AR-C338-maqbara}



\begin{quote}
Before the new one: say the Arabic for a tomb, then the Arabic for a death.
\end{quote}

\subsection*{You'll want to know: \ar{مقبرة}}
\textbf{\ar{مقبرة}} — \emph{maqbara} — \textquotedblleft{}a cemetery\textquotedblright{}.

\begin{quote}
\ar{زرنا} \ar{المقبرة} \ar{صباح} \ar{العيد}.
\end{quote}

\begin{quote}
\emph{zurnā al-maqbara ṣabāḥ al-ʿīd.}
\end{quote}

\begin{quote}
We visited the cemetery on the morning of the holiday.
\end{quote}

Plural: \ar{مقابر}, \emph{maqābir}.

\begin{cousinweb}
\textbf{\ar{مقبرة}} (\emph{maqbara}) is built on the native Arabic root \emph{q-b-r} (Lane), \textquotedblleft{}to bury\textquotedblright{}. The \emph{mafʿala} pattern names a place: \textquotedblleft{}the place of graves\textquotedblright{}.
\end{cousinweb}

\begin{grammarlens}[title={What you've built}]
One more word for telling what happened.
\end{grammarlens}

\subsection*{Your turn}
\begin{itemize}
\raggedright
  \item \emph{Say it:} \emph{maqbara}
  \item \emph{Say it:} \emph{maqbara}, once more
  \item \emph{Say it:} \emph{zurnā al-maqbara ṣabāḥ al-ʿīd.} — We visited the cemetery on the morning of the holiday.
  \item \emph{Recall:} say the Arabic for a tomb, then the Arabic for a death, then say \emph{maqbara} again
\end{itemize}

\begin{tcolorbox}[breakable,colback=teal!4,colframe=teal!35!black,title={Before you move on}]
What does \ar{مقبرة} mean? (\textquotedblleft{}A cemetery\textquotedblright{}.) Say it once more.
\end{tcolorbox}
\section[\texorpdfstring{\ar{نعش} (naʿsh)}{نعش (naʿsh)}]{\ar{نعش} (naʿsh) — a coffin, a bier}

\label{lesson:AR-C338-nash}



\begin{quote}
Before the new one: say the Arabic for a death, then the Arabic for a cemetery.
\end{quote}

\subsection*{You'll want to know: \ar{نعش}}
\textbf{\ar{نعش}} — \emph{naʿsh} — \textquotedblleft{}a coffin, a bier\textquotedblright{}.

\begin{quote}
\ar{حمل} \ar{أربعة} \ar{رجال} \ar{النعش}.
\end{quote}

\begin{quote}
\emph{ḥamala arbaʿat rijāl an-naʿsh.}
\end{quote}

\begin{quote}
Four men carried the coffin.
\end{quote}

\begin{cousinweb}
\textbf{\ar{نعش}} (\emph{naʿsh}) is built on the native Arabic root \emph{n-ʿ-sh} (Lane), \textquotedblleft{}to raise up\textquotedblright{}. The stars of the Great Bear are \emph{banāt naʿsh}, \textquotedblleft{}the daughters of the bier\textquotedblright{}.
\end{cousinweb}

\begin{grammarlens}[title={What you've built}]
One more word for telling what happened.
\end{grammarlens}

\subsection*{Your turn}
\begin{itemize}
\raggedright
  \item \emph{Say it:} \emph{naʿsh}
  \item \emph{Say it:} \emph{naʿsh}, once more
  \item \emph{Say it:} \emph{ḥamala arbaʿat rijāl an-naʿsh.} — Four men carried the coffin.
  \item \emph{Recall:} say the Arabic for a death, then the Arabic for a cemetery, then say \emph{naʿsh} again
\end{itemize}

\begin{tcolorbox}[breakable,colback=teal!4,colframe=teal!35!black,title={Before you move on}]
What does \ar{نعش} mean? (\textquotedblleft{}A coffin, a bier\textquotedblright{}.) Say it once more.
\end{tcolorbox}
